\documentclass[11pt]{scrartcl}


\usepackage[sexy]{evan}
\usepackage{import}
\usepackage[table]{xcolor}
\usepackage{xfp}

\newcommand{\gad}{\textcolor{yellow}{$\bigstar$}}
\newcommand{\bad}{\textcolor{red}{$\bigstar$}}

\definecolor{dcol0}{HTML}{C8E6C9}
\definecolor{dcol1}{HTML}{D4E9B3}
\definecolor{dcol2}{HTML}{E5ED9A}
\definecolor{dcol3}{HTML}{FFF59D}
\definecolor{dcol4}{HTML}{FFE082}
\definecolor{dcol5}{HTML}{FFCC80}
\definecolor{dcol6}{HTML}{FFAB91}
\definecolor{dcol7}{HTML}{F49890}
\definecolor{dcol8}{HTML}{E57373}
\definecolor{dcol9}{HTML}{D32F2F}

\makeatletter
\newcommand{\getcolorname}[1]{dcol#1}
\makeatother

\newcommand{\dif}[1]{%
    \edef\colorindex{\number\fpeval{floor(#1)}}%
    \edef\fulltext{#1}%
    \colorbox{\getcolorname{\colorindex}}{%
        \ifnum\colorindex>8
            \textbf{\textcolor{white}{\,\fulltext\,}}%
        \else
            \textbf{\textcolor{black}{\,\fulltext\,}}%
        \fi
    }%
}
% Variable para dificultad (inicial 0)
\newcommand{\thmdifficulty}{0}

% Comando para asignar dificultad antes del problema
\newcommand{\problemdiff}[1]{\renewcommand{\thmdifficulty}{#1}}

% Estilo del problema que incluye dificultad antes del título
\declaretheoremstyle[
    headfont=\color{blue!40!black}\normalfont\bfseries,
    headformat={%
      \dif{\thmdifficulty}\quad \NAME~\NUMBER\ifx\relax\EMPTY\relax\else\ \NOTE\fi
    },
    postheadspace=1em,
    spaceabove=8pt,
    spacebelow=8pt,
    bodyfont=\normalfont
]{problemstyle}

    \declaretheorem[style=problemstyle,name=Problema,sibling=theorem]{problema}
    \declaretheorem[style=problemstyle,name=Problema,numbered=no]{problema*}

\definecolor{yellow4}{RGB}{255, 206, 0}

\title {Pequeño Teorema de Fermat}
\subtitle{Entrenamientos Nacionales Platas y EGMOs}

\author{Emmanuel Buenrostro}


\begin{document}
\maketitle

\section{Teoria}

\begin{lemma}
    Si tienes $a,b$ enteros primos relativos, entonces $a,2a,3a, \ldots, (b-1)a$ son todos residuos distintos modulo $b$.
\end{lemma}


\begin{theorem}[Pequeño Teorema de Fermat]
 Si $p$ es primo y $a\in \mathbb{Z}$ y $(a,p)=1$ entonces 
$$a^{p-1} \equiv 1 \text{ mod } p$$
\end{theorem}

\begin{proof}
Notemos que para $a$ tal que $p|a$ si cumple porque $a^p\equiv 0 \equiv a$ mod $p$.

Para los $a$ con $(p,a)=1$. se tiene que 
$$\{a,2a,3a, \ldots, (p-1)a\}$$
es una permutación de 
$$\{1,2,3,\ldots p-1\}$$
en mod $p$. 

\underline{\textbf{Prueba.}} Todos los números $a,2a,3a,\ldots, (p-1)a$ tiene distintos modulos $p$, ya que si hay dos iguales $ia$ y $ja$ con $i\neq j$ entonces 
$$ia-ja= a(i-j) \equiv 0 \text{ mod } p$$
y como $(a,p)=1$ entonces $i-j \equiv 0$ mod $p \Rightarrow i\equiv j$ mod $p$, pero $1\leq i,j \leq p-1$ entonces $i=j$, una contradicción.
Entonces si es una permutación.  \\

Asi que si multiplicamos todos se tiene que 

$$a\cdot 2a\cdot 3a \cdots (p-1)a \equiv 1\cdot 2 \cdots (p-1) \text{ mod } p$$
entonces
$$ \Rightarrow(p-1)!a^{p-1}\equiv (p-1)! \text{ mod } p $$
$$ \Rightarrow a^{p-1}\equiv 1 \text{ mod }p$$
\end{proof}
\begin{remark*}
    Un entero $n$ es llamado \textit{Carmichael} si para todo $a$ coprimo con $n$ entonces $a^{n-1} \equiv 1 \pmod n $.
\end{remark*}

\begin{theorem}[Teorema de Euler]
 Si $a$ es primo relativo con $n$ entonces $a^{\phi(n)}\equiv 1$ mod $n$
\end{theorem}

\underline{\textbf{Prueba parecida a la 2da de Fermat.}}
Sea $a$ un entero tal que $(a,n)=1$.

Vamos a considerar el conjunto $S=\{ja: 1\leq j \leq n \text{ y } (j,n)=1\}$ y vamos a probar que es una permutación de los primos relativos a $n$.
\begin{itemize}
\item Los números $ja$ pertenicientes a $S$ son primos relativos a $n$, ya que es la multiplicacion de dos primos relativos a $n$, y entonces su multiplicación no comparte ningun factor con $n$.
\item Hay $\phi (n)$ distintos modulos primos relativos a $n$ por definición.
\item Si hay dos $1\leq j_1,j_2 \leq n$ distintos primos relativos a $n$ tal que $j_1a\equiv j_2 a$ mod $n$ entonces $(j_1-j_2)a\equiv 0$ mod $n$ entonces $n|(j_1-j_2)a$ pero como $(a,n)=1$ entonces $n|j_1-j_2$ y $j_1\equiv j_2$ mod $n$ pero como $1\leq j_1,j_2 \leq n$ entonces $j_1=j_2$ una contradicción.
\item Entonces los $\phi (n)$ $j$ son distintos y entonces hay $\phi (n)$ $ja$ diferentes y son primos relativos con $n$ entonces $S$ si es una permutacion de los primos relativos con $n$ en mod $n$.
\end{itemize}
Entonces 
$$\prod_{(j,n)=1}^{n} ja \equiv \prod_{(j,n)=1}^{n} j \text{ mod } n $$
$$\Rightarrow \prod_{(j,n)=1}^{n} ja = (\prod_{(j,n)=1}^{n}j) a^{\phi (n)} \equiv \prod_{(j,n)=1}^{n} j \text{ mod } n$$
y como $\prod_{(j,n)=1}^{n} j$ es la multiplicación de primos relativos con $n$ entonces es primo relativo con $n$ y 
$$a^{\phi(n)}\equiv 1 \text{ mod } n$$



\begin{example}
    [Caso particular de Fermat's Christmas theorem]
    Demuestra que si $p \mid a^2+1$ para algun entero $a$, entonces $p=2$ o $p \equiv 1 \pmod 4$. 
\end{example}

\begin{soln}
La divisibilidad nos da que $a^2 \equiv -1 \pmod p$. \\
Si $p=4m+3$ entonces 
\[a^{p-1}\equiv  (a^2)^{2m+1}\equiv (-1)^{2m+1} \equiv -1 \pmod p\]
Lo cual es una contradicción con el Pequeño Teorema de Fermat
\end{soln}
\begin{remark*}
Sugerencia de Eli para ponernos modo navideño
\end{remark*}


Si tienes una expresion del estilo $a^n+n$ o similares puedes hacer que sea $0 \pmod p $ tomando $n \equiv 1 \pmod{p-1}$ y $n \equiv -a \pmod p$.  



\newpage 
\section{Problemas}
\epigraph{“You never know how strong you are until being strong is the only choice you have.”}{Libro de Justin Stevens}
    
\subsection{Nivel I}

\problemdiff{1}
\begin{problem}
Calcula $2^{1000} \pmod {997}$.
\end{problem}

\problemdiff{1}

\begin{problem}
Encuentra todos los $x$ tales que $x^{35}+5x^{19}+11x^{3}$ es divisible entre 17.
\end{problem}

\problemdiff{1}
\begin{problem}
Encuentra todos los primos $p$ para los cuales $13^{2p-1}+17$ es divisible entre $p$.
\end{problem}

\problemdiff{2}
\begin{problem}
Sea $f(x_1,x_2,\ldots, x_n)$ un polinomio con coeficientes enteros. Prueba que 
$$f(x_1,x_2, \ldots ,x_n)^p \equiv f(x_1^p,x_2^p, \ldots, x_n^p) \pmod p$$
para $p$ primo.
\end{problem}

\problemdiff{2}
\begin{problem}
   % [PuMAC]
    Calcula los ultimos 3 digitos de $2008^{2007^{2006^{\ldots^{2^1}}}}$
\end{problem}


\problemdiff{3}
\begin{problem}
  %  [nhhlqd]
    Encuentra todos los pares de enteros positivos $(p,n)$ con $p$ primo tal que: 
    \[p-1 \mid n+3\]
    y 
    \[p \mid 2^{n-1}+3\]
\end{problem}

\problemdiff{3}
\begin{problem}
  %  [Rioplatense 2004/4]
    Encuentra el número de enteros $n>1$ tales que $a^{13}-a$ es divisible entre $n$ para todo entero positivo $a$. 
\end{problem}

\problemdiff{3}
\begin{problem}
  %  [Canada 1983/4]
    Prueba que para todo primo $p$, existen infinitos enteros $n$ tales que $p \mid 2^{n}-n$.
\end{problem}

\subsection{Nivel II}

\problemdiff{4}
\begin{problem}
  %  [IMO 2005/4]
    
    Encuentra todos los enteros positivos que son primos relativos con todos los términos de la secuencia infinita \[ a_n=2^n+3^n+6^n -1,\ n\geq 1. \]
    
\end{problem}


\problemdiff{4}
\begin{problem}
  %  [Libro de Justin Stevens]
    Sea $n$ un entero positivo con $n\geq 3$. Prueba que 
    \[n^{n^{n^{n}}}-n^{n^{n}}\]
    es divisible entre 1989.
\end{problem}


\problemdiff{4}
\begin{problem}
 %   [Balkan 1999/2]
    Sea $p>2$ un número primo con $3 \mid p-2$. Sea
    \[S=\{y^2-x^3-1 \mid 0 \leq x ,y \leq p-1 \cap x,y \in \mathbb{Z}\}\]
    Prueba que a lo más $p$ elementos de $S$ son divisibles por $p$.
\end{problem}

\problemdiff{4}
\begin{problem}
 %   [JBMO 2020/4]
    Encuentra todos los primos $p,q$ tal que
    \[1+\frac{p^q-q^p}{p+q}\]
    es un número primo.
\end{problem}

\problemdiff{4}
\begin{problem}
    Prueba que 561 es un número Carmichael. \\ 
    \textit{Nota: Abajo de la prueba de Fermat viene que es un Carmichael}.
\end{problem}

\subsection{Nivel III}
\problemdiff{5}
\begin{problem}
 %   [OMM 2010/6]

    Sean $p, q, r$ números primos positivos distintos. Muestra que si $pqr$ divide a
\[(pq)^r+(qr)^p+(rp)^q-1,\]
entonces $(pqr)^3$ divide a
\[3((pq)^r+(qr)^p+(rp)^q-1).\]
\end{problem}

\problemdiff{5}
\begin{problem}
 %   [Korselt's Criterion]
    Un número $n$ es Carmichael si y solo si $n=p_1p_2\cdots p_r$, donde cada $p_i$ es distinto ($n$ es libre de cuadrados) y $p_i-1 \mid n-1$ para todo $1\leq i \leq r$.
\end{problem}

\problemdiff{5}
\begin{problem}
 %   [Iran Finals 2016/ N1]
    Sean $p,q$ primos con $q$ impar. Prueba que existe un entero $x$ tal que 
    \[ q \mid (x+1)^p-x^p \] 
    si y solo si 
    \[q \equiv 1 \pmod p\]
\end{problem}


\problemdiff{5.5}
\begin{problem}
%    [IMO SL 2017/N2]
    Sea $p \geq 2$ un número primo. Adán y Cesár juegan el siguiente juego alternadamente: En cada movimiento, el jugador del turno 
    escoge un número $i$ en el conjunto $\{0,1,2,\ldots, p-1\}$ que no a sido escogido antes por ninguno de los dos jugadores y escoge un elemento $a_i$ del conjunto $\{0,1,2,3,4,5,6,7,8,9\}$. Adán juega primero. El juego termina despues de que todos los indices han sido escogidos. Entonces el siguiente número es calculado 
    $$M=a_0+a_110+a_210^2+\cdots+a_{p-1}10^{p-1}= \sum_{i=0}^{p-1}a_i.10^i$$

    La meta de Adán es hacer $M$ multiplo de $p$, y la meta de César es evitar eso. \\
    Prueba que Adán tiene la estrategia ganadora. 
\end{problem}


\subsection{Nivel IV}
\problemdiff{6}
\begin{problem}
 %   [USAMO 1991/3]
    Prueba que para cualquier entero $n\geq 1$, la sucesión 
    \[2, 2^2, 2^{2^2}, 2^{2^{2^{2}}}, \ldots \pmod n\]
    es eventualmente constante. 
\end{problem}

\problemdiff{6}
\begin{problem}
  %  [OMMFEM 2024/8]
    Encuentra todos los enteros positivos $n$ tal que entre los $n$ números 
    \[ 2n + 1, \, 2^2 n + 1, \, \ldots, \, 2^n n + 1 \]
    hay $n, n-1$ o $n-2$ primos.
\end{problem}

\problemdiff{6}
\begin{problem}
 %   [IMO SL 2024/N3]
    Determina todas las sucesiones $a_1, a_2, \ldots$ de enteros positivos tales que para cualquier pareja de enteros positivos $m \leq n$, 
    \[ \frac{a_m + a_{m+1} + \cdots + a_n}{n-m+1}\quad\text{y}\quad (a_ma_{m+1}\cdots a_n)^{\frac{1}{n-m+1}}\]
    son ambos enteros.
\end{problem}

\problemdiff{7}
\begin{problem}
 %   [IMO SL 2005/N6]
    Sean $a,b$ enteros positivos tal que $b^n +n$ es multiplo de $a^n+n$ para todos los enteros $n$. Prueba que $a=b$.
\end{problem}

\problemdiff{7}
\begin{problem}
 %   [IMO SL 2006/N5]
    Encuentra todas las soluciones enteras de 
    \[\frac {x^{7} - 1}{x - 1} = y^{5} - 1.\]
\end{problem}

\problemdiff{7}
\begin{problem}
  %  [IMO 2025/3]
    Sea $\mathbb{N}$ el conjunto de enteros positivos. Una funcion $f\colon\mathbb{N}\to\mathbb{N}$ se llama \textit{bonza}
    \[
f(a)~~\text{divide a}~~b^a-f(b)^{f(a)}
\]
para todos los enteros positivos $a,b$. \\
Determina la constante real $c$ tal que $f(n) \leq cn $ para todas las funciones bonza $f$ y enteros positivos $n$.
\end{problem}
\end{document}